% !TeX root = mainstarkov.tex
% !TeX program = xelatex
% !TeX encoding = UTF-8
\section*{ЗАКЛЮЧЕНИЕ}
\addcontentsline{toc}{section}{ЗАКЛЮЧЕНИЕ}

В рамках выпускной квалификационной работы выполнена программная реализация и тестирование алгоритмов динамической маршрутизации для системы построения эвакуационных маршрутов в замкнутых пространствах. Все поставленные задачи решены в полном объёме.

Проведён анализ существующих алгоритмов поиска кратчайших путей на графах. На основании сравнительного анализа по критериям оптимальности, производительности, поддержки динамического исключения узлов и сложности реализации выбран алгоритм Дейкстры как наилучшее решение для задачи эвакуационной маршрутизации. Алгоритм обеспечивает гарантию оптимальности при неотрицательных весах рёбер, эффективную работу на разреженных графах, характерных для планов зданий, и допускает модификацию для динамического исключения заблокированных узлов.

Разработаны структуры данных для представления навигационного графа: классы \texttt{NavGraph}, \texttt{NavNode} и \texttt{NavEdge}, обеспечивающие эффективное хранение и обработку топологии здания. Класс \texttt{NavNode} содержит полный набор атрибутов узла: координаты, номер этажа, идентификатор здания, флаги типа и группу перехода. Класс \texttt{NavEdge} хранит вес ребра с автоматическим вычислением евклидова расстояния для проходов. Реализована загрузка и сохранение графа в формате JSON средствами \texttt{System.Text.Json}.

Реализован алгоритм Дейкстры с модификациями для поддержки динамического исключения узлов. Множество исключённых узлов передаётся как параметр метода \texttt{FindPath} и проверяется при релаксации рёбер без модификации исходного графа. Конечный узел маршрута никогда не исключается. Алгоритм использует стандартные структуры данных .NET (\texttt{SortedSet}, \texttt{Dictionary}, \texttt{HashSet}) и обеспечивает время выполнения менее 200 мс для графов до 500 узлов, что соответствует требованиям реального времени. Оптимизация списка смежности и обработка устаревших записей очереди повышают эффективность реализации.

Разработан алгоритм поиска ближайшего эвакуационного выхода, учитывающий текущее положение пользователя, принадлежность к зданию и динамически изменяющиеся препятствия. Алгоритм приоритизирует выходы в текущем здании. Алгоритм интегрирован с сервисом экстренной эвакуации \texttt{EmergencyService} и ViewModel приложения IndoorNav.

Разработан механизм формирования множества исключённых узлов \texttt{BuildExcludeSet}, объединяющий заблокированные пользователем узлы и выходы, не являющиеся целевыми. Это предотвращает прохождение маршрута через другие выходы и обеспечивает корректное перестроение маршрута при блокировке участков.

Проведено комплексное тестирование алгоритмической части: модульное тестирование четырёх модулей маршрутизации (построение маршрута, поиск эвакуационного выхода, вычисление весов рёбер и формирование множества исключений), тестирование производительности на графах от 50 до 500 узлов, а также интеграционное тестирование на реальных данных графа здания университета. Результаты подтвердили корректность работы алгоритмов во всех сценариях, включая граничные случаи (пустой граф, несвязный граф, совпадение стартовой и конечной точек), и соответствие требованиям производительности.

Таким образом, в ходе выпускной квалификационной работы реализованы и верифицированы алгоритмы динамической маршрутизации для системы построения эвакуационных маршрутов в замкнутых пространствах. Полученные результаты могут быть применены в навигационных системах внутри зданий, приложениях экстренной эвакуации и иных программных комплексах, требующих эффективного построения маршрутов с учётом динамических ограничений. Дальнейшее развитие системы может включать реализацию алгоритма A* с допустимой эвристикой, поддержку двунаправленного поиска и интеграцию с серверной частью для распределённого построения маршрутов.
